% GENERATED FILE. Edit canonical lessons, then run npm run generate:books.
% canonical-source-hash: fnv1a64:b56feb0dbe5cdd75
% canonical-lessons: PA-C332-bhatija, PA-C332-bhanja, PA-C332-guandhi, PA-C332-saheli, PA-C332-pukarna

\chapter{Nephews, Friends and Neighbours}
\label{ch:pa-a2-nephews-friends-and-neighbours}
\hlchaptermodality{332}

\begin{chapteropening}
\textbf{By the end of this chapter:} \emph{I can name a nephew, a friend and the people next door.}

Complete the last lesson of chapter 332: I can name a nephew, a friend and the people next door.
\end{chapteropening}

\section[\texorpdfstring{\pa{ਭਤੀਜਾ} (bhatījā)}{ਭਤੀਜਾ (bhatījā)}]{\pa{ਭਤੀਜਾ} (bhatījā) — a brother's son}

\label{lesson:PA-C332-bhatija}



\begin{quote}
Before the new one: say the Punjabi for a mother's sister, then the Punjabi for a mother's sister's husband.
\end{quote}

\subsection*{You'll want to know: \pa{ਭਤੀਜਾ}}
\textbf{\pa{ਭਤੀਜਾ}} — \emph{bhatījā} — \textquotedblleft{}a brother's son\textquotedblright{}.

\begin{quote}
\pa{ਮੇਰਾ} \pa{ਭਤੀਜਾ} \pa{ਹਰ} \pa{ਰੋਜ਼} \pa{ਸਕੂਲ} \pa{ਜਾਂਦਾ} \pa{ਹੈ।}
\end{quote}

\begin{quote}
\emph{merā bhatījā har roz sakūl jāndā hai.}
\end{quote}

\begin{quote}
My brother's son goes to school every day.
\end{quote}

A masculine noun.

\begin{cousinweb}
\textbf{\pa{ਭਤੀਜਾ}} (\emph{bhatījā}) is inherited from Sanskrit \emph{bhrātrīya}, \textquotedblleft{}a brother's son\textquotedblright{}; CDIAL lists Punjabi \emph{bhatījā}, \textquotedblleft{}brother's son\textquotedblright{}, under it (CDIAL 9672).
\end{cousinweb}

\begin{grammarlens}[title={What you've built}]
One more word for this chapter.
\end{grammarlens}

\subsection*{Your turn}
\begin{itemize}
\raggedright
  \item \emph{Say it:} \emph{bhatījā}
  \item \emph{Say it:} \emph{bhatījā}, once more
  \item \emph{Say it:} \emph{merā bhatījā har roz sakūl jāndā hai.} — My brother's son goes to school every day.
  \item \emph{Recall:} say the Punjabi for a mother's sister, then the Punjabi for a mother's sister's husband, then say \emph{bhatījā} again
\end{itemize}

\begin{tcolorbox}[breakable,colback=teal!4,colframe=teal!35!black,title={Before you move on}]
What does \pa{ਭਤੀਜਾ} mean? (\textquotedblleft{}A brother's son\textquotedblright{}.) Say it once more.
\end{tcolorbox}
\section[\texorpdfstring{\pa{ਭਾਣਜਾ} (bhāṇjā)}{ਭਾਣਜਾ (bhāṇjā)}]{\pa{ਭਾਣਜਾ} (bhāṇjā) — a sister's son}

\label{lesson:PA-C332-bhanja}



\begin{quote}
Before the new one: say the Punjabi for a mother's sister's husband, then the Punjabi for a brother's son.
\end{quote}

\subsection*{You'll want to know: \pa{ਭਾਣਜਾ}}
\textbf{\pa{ਭਾਣਜਾ}} — \emph{bhāṇjā} — \textquotedblleft{}a sister's son\textquotedblright{}.

\begin{quote}
\pa{ਮੇਰਾ} \pa{ਭਾਣਜਾ} \pa{ਕੱਲ੍ਹ} \pa{ਸਾਡੇ} \pa{ਘਰ} \pa{ਆਇਆ।}
\end{quote}

\begin{quote}
\emph{merā bhāṇjā kall sāḍe ghar āiā.}
\end{quote}

\begin{quote}
My sister's son came to our house yesterday.
\end{quote}

A masculine noun.

\begin{cousinweb}
\textbf{\pa{ਭਾਣਜਾ}} (\emph{bhāṇjā}) is inherited from Sanskrit \emph{bhāgineya}, \textquotedblleft{}a sister's son\textquotedblright{}; CDIAL lists Punjabi \emph{bhāṇjā}, \textquotedblleft{}sister's son\textquotedblright{}, under it (CDIAL 9433).
\end{cousinweb}

\begin{grammarlens}[title={What you've built}]
One more word for this chapter.
\end{grammarlens}

\subsection*{Your turn}
\begin{itemize}
\raggedright
  \item \emph{Say it:} \emph{bhāṇjā}
  \item \emph{Say it:} \emph{bhāṇjā}, once more
  \item \emph{Say it:} \emph{merā bhāṇjā kall sāḍe ghar āiā.} — My sister's son came to our house yesterday.
  \item \emph{Recall:} say the Punjabi for a mother's sister's husband, then the Punjabi for a brother's son, then say \emph{bhāṇjā} again
  \item \emph{Write it:} copy \textbf{\pa{ਭਾਣਜਾ}} once, checking the whole word
\end{itemize}

\begin{tcolorbox}[breakable,colback=teal!4,colframe=teal!35!black,title={Before you move on}]
What does \pa{ਭਾਣਜਾ} mean? (\textquotedblleft{}A sister's son\textquotedblright{}.) Say it once more.
\end{tcolorbox}
\section[\texorpdfstring{\pa{ਗੁਆਂਢੀ} (guāṇḍhī)}{ਗੁਆਂਢੀ (guāṇḍhī)}]{\pa{ਗੁਆਂਢੀ} (guāṇḍhī) — a neighbour}

\label{lesson:PA-C332-guandhi}



\begin{quote}
Before the new one: say the Punjabi for a brother's son, then the Punjabi for a sister's son.
\end{quote}

\subsection*{You'll want to know: \pa{ਗੁਆਂਢੀ}}
\textbf{\pa{ਗੁਆਂਢੀ}} — \emph{guāṇḍhī} — \textquotedblleft{}a neighbour\textquotedblright{}.

\begin{quote}
\pa{ਸਾਡੇ} \pa{ਗੁਆਂਢੀ} \pa{ਬਹੁਤ} \pa{ਚੰਗੇ} \pa{ਲੋਕ} \pa{ਹਨ।}
\end{quote}

\begin{quote}
\emph{sāḍe guāṇḍhī bahut chaṅge lok han.}
\end{quote}

\begin{quote}
Our neighbours are very good people.
\end{quote}

A masculine or feminine noun, by the person meant.

\begin{cousinweb}
\textbf{\pa{ਗੁਆਂਢੀ}} (\emph{guāṇḍhī}) goes back to a reconstructed \emph{*grāmārdha}, \textquotedblleft{}the vicinity of a village\textquotedblright{}, built from \emph{grāma}, \textquotedblleft{}village\textquotedblright{}, and \emph{ardha}, \textquotedblleft{}half, side\textquotedblright{}; CDIAL lists Punjabi \emph{gavā̃ḍhī}, \textquotedblleft{}neighbour\textquotedblright{}, under it (CDIAL 4376).
\end{cousinweb}

\begin{grammarlens}[title={What you've built}]
One more word for this chapter.
\end{grammarlens}

\subsection*{Your turn}
\begin{itemize}
\raggedright
  \item \emph{Say it:} \emph{guāṇḍhī}
  \item \emph{Say it:} \emph{guāṇḍhī}, once more
  \item \emph{Say it:} \emph{sāḍe guāṇḍhī bahut chaṅge lok han.} — Our neighbours are very good people.
  \item \emph{Recall:} say the Punjabi for a brother's son, then the Punjabi for a sister's son, then say \emph{guāṇḍhī} again
  \item \emph{Write it:} copy \textbf{\pa{ਗੁਆਂਢੀ}} once, checking the whole word
\end{itemize}

\begin{tcolorbox}[breakable,colback=teal!4,colframe=teal!35!black,title={Before you move on}]
What does \pa{ਗੁਆਂਢੀ} mean? (\textquotedblleft{}A neighbour\textquotedblright{}.) Say it once more.
\end{tcolorbox}
\section[\texorpdfstring{\pa{ਸਹੇਲੀ} (sahelī)}{ਸਹੇਲੀ (sahelī)}]{\pa{ਸਹੇਲੀ} (sahelī) — a woman's female friend}

\label{lesson:PA-C332-saheli}



\begin{quote}
Before the new one: say the Punjabi for a sister's son, then the Punjabi for a neighbour.
\end{quote}

\subsection*{You'll want to know: \pa{ਸਹੇਲੀ}}
\textbf{\pa{ਸਹੇਲੀ}} — \emph{sahelī} — \textquotedblleft{}a woman's female friend\textquotedblright{}.

\begin{quote}
\pa{ਮੇਰੀ} \pa{ਭੈਣ} \pa{ਆਪਣੀ} \pa{ਸਹੇਲੀ} \pa{ਨਾਲ} \pa{ਬਾਜ਼ਾਰ} \pa{ਗਈ।}
\end{quote}

\begin{quote}
\emph{merī bhaiṇ āpṇī sahelī nāl bāzār gaī.}
\end{quote}

\begin{quote}
My sister went to the market with her friend.
\end{quote}

A feminine noun.

\begin{cousinweb}
\textbf{\pa{ਸਹੇਲੀ}} (\emph{sahelī}) is inherited from Sanskrit \emph{sakhī}, \textquotedblleft{}a woman's confidante\textquotedblright{}; CDIAL lists Punjabi \emph{sahelī}, \textquotedblleft{}woman's female friend\textquotedblright{}, under it (CDIAL 13074).
\end{cousinweb}

\begin{grammarlens}[title={What you've built}]
One more word for this chapter.
\end{grammarlens}

\subsection*{Your turn}
\begin{itemize}
\raggedright
  \item \emph{Say it:} \emph{sahelī}
  \item \emph{Say it:} \emph{sahelī}, once more
  \item \emph{Say it:} \emph{merī bhaiṇ āpṇī sahelī nāl bāzār gaī.} — My sister went to the market with her friend.
  \item \emph{Recall:} say the Punjabi for a sister's son, then the Punjabi for a neighbour, then say \emph{sahelī} again
  \item \emph{Write it:} copy \textbf{\pa{ਸਹੇਲੀ}} once, checking the whole word
\end{itemize}

\begin{tcolorbox}[breakable,colback=teal!4,colframe=teal!35!black,title={Before you move on}]
What does \pa{ਸਹੇਲੀ} mean? (\textquotedblleft{}A woman's female friend\textquotedblright{}.) Say it once more.
\end{tcolorbox}
\section[\texorpdfstring{\pa{ਪੁਕਾਰਨਾ} (pukārnā)}{ਪੁਕਾਰਨਾ (pukārnā)}]{\pa{ਪੁਕਾਰਨਾ} (pukārnā) — to call out to someone}

\label{lesson:PA-C332-pukarna}



\begin{quote}
Before the new one: say the Punjabi for a neighbour, then the Punjabi for a woman's female friend.
\end{quote}

\subsection*{You'll want to know: \pa{ਪੁਕਾਰਨਾ}}
\textbf{\pa{ਪੁਕਾਰਨਾ}} — \emph{pukārnā} — \textquotedblleft{}to call out to someone\textquotedblright{}.

\begin{quote}
\pa{ਮੇਰਾ} \pa{ਭਤੀਜਾ} \pa{ਮੈਨੂੰ} \pa{ਗਲੀ} \pa{ਤੋਂ} \pa{ਪੁਕਾਰਦਾ} \pa{ਹੈ।}
\end{quote}

\begin{quote}
\emph{merā bhatījā mainū̃ galī ton pukārdā hai.}
\end{quote}

\begin{quote}
My brother's son calls out to me from the lane.
\end{quote}

Past (I …): \pa{ਮੈਂ} \pa{ਗੁਆਂਢੀ} \pa{ਨੂੰ} \pa{ਪੁਕਾਰਿਆ}, \emph{maiṁ guāṇḍhī nū̃ pukāriā}.

\begin{cousinweb}
\textbf{\pa{ਪੁਕਾਰਨਾ}} (\emph{pukārnā}) is inherited from a reconstructed \emph{*pukkār}, \textquotedblleft{}a shout\textquotedblright{}; CDIAL lists Punjabi \emph{pukārnā}, \textquotedblleft{}shout\textquotedblright{}, under it (CDIAL 8246).
\end{cousinweb}

\begin{grammarlens}[title={What you've built}]
One more word for this chapter.
\end{grammarlens}

\subsection*{Your turn}
\begin{itemize}
\raggedright
  \item \emph{Say it:} \emph{pukārnā}
  \item \emph{Say it:} \emph{pukārnā}, once more
  \item \emph{Say it:} \emph{merā bhatījā mainū̃ galī ton pukārdā hai.} — My brother's son calls out to me from the lane.
  \item \emph{Recall:} say the Punjabi for a neighbour, then the Punjabi for a woman's female friend, then say \emph{pukārnā} again
  \item \emph{Write it:} copy \textbf{\pa{ਪੁਕਾਰਨਾ}} once, checking the whole word
\end{itemize}

\begin{tcolorbox}[breakable,colback=teal!4,colframe=teal!35!black,title={Before you move on}]
What does \pa{ਪੁਕਾਰਨਾ} mean? (\textquotedblleft{}To call out to someone\textquotedblright{}.) Say it once more.
\end{tcolorbox}
